% answers_short.tex —— 2008 年·普通高等学校招生考试·江苏卷 · 简短答案速查
% 由 tools/make_short_answers.py 自 answers.tex 的 \gksolhead/\gkanswer 提取，
% 并与 2008/2008_数学真题_江苏_答案要点.md 逐题比对后生成，勿手改（改源 MD 后重跑）。

\begin{gknotice}
\textbf{使用说明}：本节只给最终答案便于快速核对，完整推导见解析版。答案为编写者独立推导，\textbf{非官方标准答案}。
\end{gknotice}

\gkpart{填空题}
\begin{itemize}[leftmargin=2.6em]
  \item[\textbf{1.}] $10$。
  \item[\textbf{2.}] $\dfrac1{12}$。
  \item[\textbf{3.}] $1$。
  \item[\textbf{4.}] $6$。
  \item[\textbf{5.}] $7$。
  \item[\textbf{6.}] $\dfrac{\pi}{16}$。
  \item[\textbf{7.}] $6.42$。
  \item[\textbf{8.}] $\ln2-1$。
  \item[\textbf{9.}] $\dfrac1c-\dfrac1b$。
  \item[\textbf{10.}] $\dfrac{(n-1)n}2+3$。
  \item[\textbf{11.}] $3$。
  \item[\textbf{12.}] $\dfrac{\sqrt2}2$。
  \item[\textbf{13.}] $2\sqrt2$。
  \item[\textbf{14.}] $4$。
\end{itemize}

\gkpart{解答题}
\begin{itemize}[leftmargin=2.6em]
  \item[\textbf{15.}] （1）$-3$；（2）$\dfrac{3\pi}4$。
  \item[\textbf{16.}] （1）（2）证明见解析。
  \item[\textbf{17.}] （1）① $y=\dfrac{20}{\cos\theta}-10\tan\theta+10$（$0\leqslant\theta\leqslant\frac{\pi}4$）；② $y=2\sqrt{100+(10-x)^2}+x$（$0\leqslant x\leqslant10$）；（2）$OA=OB=PB=PA$ 即 $O$ 在 $AB$ 中垂线上距 $AB$ $\dfrac{10\sqrt3}3$ km 处，最短总长度 $\left(10+10\sqrt3\right)$ km。
  \item[\textbf{18.}] （1）$b<1$ 且 $b\neq0$；（2）$x^2+y^2+2x-(b+1)y+b=0$；（3）过定点 $(0,1)$ 与 $(-2,1)$。
  \item[\textbf{19.}] （1）① $-4$ 或 $1$；② $n=4$；（2）证明见解析。
  \item[\textbf{20.}] （1）$|p_1-p_2|\leqslant\log_32$；（2）证明见解析。
\end{itemize}

\gkpart{附加题}
\begin{itemize}[leftmargin=2.6em]
  \item[\textbf{21.}] A 证明见解析；B $x^2+y^2=1$；C $2$；D 证明见解析。
  \item[\textbf{22.}] $\lambda\in\left(\dfrac13,1\right)$。
  \item[\textbf{23.}] （1）（2）证明见解析。
\end{itemize}
